% ============================================================
% Plantilla: Artículo con el Manual de Estilo de Chicago (autor-fecha)
%   \autocite{clave}         -> (Autor año)
%   \autocite[45]{clave}     -> (Autor año, 45)
%   \textcite{clave}         -> Autor (año)
% Recompilar ejecuta biber automáticamente.
% ============================================================
\documentclass[12pt,a4paper]{article}

\usepackage[utf8]{inputenc}
\usepackage[T1]{fontenc}
\usepackage{mathptmx}
\usepackage[spanish]{babel}
\usepackage[margin=2.54cm]{geometry}
\usepackage{setspace}
\usepackage{csquotes}
\usepackage[authordate,backend=biber]{biblatex-chicago}
\usepackage[hidelinks]{hyperref}

\addbibresource{referencias.bib}
\doublespacing

\title{Título del artículo}
\author{Nombre del Autor\\\small Universidad}
\date{\today}

\begin{document}

\maketitle

\section{Introducción}
El sistema autor-fecha de Chicago es habitual en ciencias sociales
\autocite{chicago2017}. Para citar una página concreta se indica después del año
\autocite[45]{lamport1994}. También puede citarse en la narración: según
\textcite{knuth1984}, la tipografía influye en la lectura.

\section{Desarrollo}
Desarrolla los argumentos principales con el respaldo de las fuentes
\autocite{goossens2004}.

\section{Conclusiones}
Resume las ideas principales.

\printbibliography[title=Referencias]

\end{document}
